\section{Preliminaries}
\label{sec:prelim}

\subsection{Recurrent world models}
\label{sec:rssm}

Our world model is the recurrent state space model (RSSM)~\cite{hafner2019planet} of DreamerV3~\cite{hafner2025dreamerv3}.
An RSSM learns the dynamics of an environment in a compact latent space, and the agent learns its policy from trajectories imagined in this space; with one fixed configuration, DreamerV3 learns more than 150 tasks, from continuous control to Atari and Minecraft~\cite{hafner2025dreamerv3}.
At environment step $t$, a recurrent core updates the deterministic state $h_t$, from which the posterior and the prior infer the stochastic latent $z_t$:
\begin{equation}
\begin{aligned}
h_t &= f(h_{t-1}, z_{t-1}, a_{t-1}),\\
z_t &\sim q(z_t \mid h_t, e_t), \quad \hat z_t \sim p(\hat z_t \mid h_t),
\end{aligned}
\label{eq:rssm}
\end{equation}
where $e_t$ is the embedding of the observation $o_t$, and heads on $(h_t,z_t)$ predict the observation, the reward and the continuation of the episode.
The core $f$, a GRU in DreamerV3, is the only path from one step to the next, so in a partially observable task $h_t$ is the memory of the agent, both when the model filters observations and when it imagines.

\subsection{Spiking neurons and spiking world models}
\label{sec:swm}

SNNs compute with neurons that integrate synaptic input into a membrane potential and communicate by binary spikes~\cite{maass1997networks}.
The leaky integrate-and-fire (LIF) neuron, discretized with the Euler method, updates its potential $v$ in simulation step $n$ as
\begin{equation}
v[n] = \lambda\, v[n-1] + (1-\lambda)\, I[n], \quad s[n] = \Hv\big(v[n] - \vartheta\big),
\label{eq:lif}
\end{equation}
where $I[n]$ is the synaptic current of the input spikes, $\lambda=1-1/\tau$ is the leak of the time constant $\tau$, $\vartheta$ is the threshold and $\Hv$ is the Heaviside step, trained with a surrogate derivative~\cite{neftci2019surrogate,wu2018stbp}.
A spike resets the potential, $v[n]\leftarrow(1-s[n])\,v[n]$.

\swm~\cite{sun2025spikingwm} builds the RSSM from LIF layers that run $T=5$ simulation steps per environment step, and replaces the GRU core with multicompartment neurons (MCNs), so that $h_t\in\{0,1\}^{T\times N}$ is the spike output of the core.
In an MCN, a basal and an apical dendrite integrate synaptic current as in \cref{eq:lif}, and the soma integrates the basal potential under a gate set by the apical one:
\begin{align}
v^{c}[n] &= \lambda_c\, v^{c}[n-1] + (1-\lambda_c)\, I^{c}[n], \label{eq:mcn}\\
u[n] &= (1-\gamma[n])\, u[n-1] + \gamma[n]\big(v^{\mathrm{bas}}[n] + I^{\mathrm{som}}[n]\big), \nonumber
\end{align}
where $c\in\{\mathrm{bas},\mathrm{api}\}$, $I^{c}[n] = W_c\,[\,x_t[n];\,h_{t-1}[n]\,]$ combines the input spikes $x_t[n]$, which encode $(z_{t-1},a_{t-1})$, with the spikes of the core at the previous environment step, $I^{\mathrm{som}}[n]=W_s x_t[n]$, and $\gamma[n]=\sigma(v^{\mathrm{api}}[n])/\tau_s$.
The soma fires $s[n]=\Hv(u[n]-\vartheta)$.
Two resets act on this neuron: the step reset sets all compartments to zero at the start of every environment step, and the spike reset clears the soma and both dendrites whenever the soma fires.\footnote{The equations of~\cite{sun2025spikingwm} show only the somatic reset; the dendritic reset is part of its official implementation (\texttt{MCNode.calc\_spike}), which all our \swm runs use.}

\subsection{Reset amnesia}
\label{sec:amnesia}

A spiking core passes information from one environment step to the next through the spikes $h_{t-1}$ it feeds back and through the states of its neurons, and only the neuron states can hold a graded memory over many steps.
We call the loss of this memory caused by the resets and the write gain of existing spiking neurons reset amnesia.

The step reset, used by spiking agents~\cite{tang2020popsan,liu2023dsqn,sun2025spikingwm}, removes the neuron states between environment steps.
The core then reduces to a binary recurrent network $h_t=\Phi(x_t,h_{t-1})$ that must encode everything it knows about the past again into the $T\times N$ bits of $h_{t-1}$ at every step, and the gradient from a reward $L$ steps after a cue crosses $L$ surrogate derivatives.
The time constants act only within the $T$ simulation steps of one environment step, which is why changing them barely changes the performance of \swm~\cite{sun2025spikingwm}.

The spike reset of the LIF neuron and of the MCN truncates a trace at the next spike.
For a compartment with this reset, the influence of a past input on the potential is
\begin{equation}
\frac{\partial v[n]}{\partial I[n-j]} = (1-\lambda)\,\lambda^{j}\prod_{i=1}^{j}\big(1-s[n-i]\big),
\label{eq:spike_reset}
\end{equation}
which vanishes once the neuron has fired, so the memory of the compartment lasts at most until its next spike.
Even without resets, the Euler update writes the input with the gain $1-\lambda=1/\tau$: for an input of unit variance that is uncorrelated over time, the state has the stationary variance $(1-\lambda)/(1+\lambda)\approx1/(2\tau)$, so a compartment with a long time constant keeps a cue but stores it too faintly to drive spikes.

Reset amnesia thus gives a design rule: the memory of a spiking world model must reside in a state that persists across environment steps, is not reset by spikes, and is written with a gain that does not vanish for long time constants.
\Cref{sec:causal} tests each mechanism with interventions on \swm.
